\documentclass[11pt]{article}
\usepackage[margin=1in]{geometry}
\usepackage{amsmath,amssymb,amsthm,booktabs}
\usepackage[hidelinks]{hyperref}
\usepackage{xurl}
\setlength{\emergencystretch}{3em}
\newtheorem{theorem}{Theorem}
\newtheorem{lemma}[theorem]{Lemma}
\newtheorem{corollary}[theorem]{Corollary}
\newcommand{\C}{\mathbb C}
\newcommand{\GL}{\operatorname{GL}}
\newcommand{\Ad}{\operatorname{Ad}}
\newcommand{\ar}{\operatorname{ar}_{\C}}
\newcommand{\rr}{r_{\mathrm{refl}}}
\newcommand{\rank}{\operatorname{rank}}
\title{Finite complex reflection-class bounds are dominated by augmentation rank}
\author{Written comparison and primary classification import}
\date{17 November 2026}
\begin{document}
\maketitle

Let $\ar(K)$ denote ordinary $U=1$ augmentation rank over $\C$,
with the meridian value different from one and arbitrary nonzero
longitude value, as in Cornwell's Definition~1.8 \cite{cor17}.
For a conjugacy class $C$ generating a finite group $W$, let
$d_C(W)$ be the minimum number of members of $C$ generating $W$.

\begin{theorem}\label{thm:domination}
Let $K$ be any knot and let
$\rho:\pi_1(S^3\setminus K)\to\GL(V)$ be a finite-dimensional complex
representation with finite image $W$. Suppose a meridian $m$ has
$\rank(\rho(m)-I)=1$, and let $C$ be its conjugacy class in $W$.
Then
\[
                         d_C(W)\le \ar(K).
\]
Consequently a meridional lower bound from one finite complex
reflection-class image is never larger than ordinary complex
augmentation rank.
\end{theorem}

Here a complex reflection is a nonidentity finite-order linear
transformation fixing a hyperplane pointwise. The hypothesis is
about one finite complex image whose meridian has displacement rank
one. The conclusion does not compare all finite-group or complex
linear methods with augmentation rank. Product or joint-kernel
targets with combined meridian displacement of larger rank, and
positive-characteristic representations, are outside its scope.

\section{The essential representation and the class condition}

\begin{lemma}\label{lem:active}
Let $G$ be normally generated by $m$, and let
$\rho:G\to\GL(V)$ have finite complex image $W$ with
$\rank(\rho(m)-I)=1$. Then there is exactly one nontrivial
irreducible summand $U$. It is faithful on $W$, and every member of
the conjugacy class $C$ of $\rho(m)$ is a reflection on $U$ with one
common nontrivial eigenvalue $\mu\ne1$. The class $C$ generates $W$.
\end{lemma}
\begin{proof}
Average a positive Hermitian form over $W$ and decompose into
irreducible summands. The ranks of $\rho(m)-I$ on these summands
add to one, so exactly one summand is active. On every other
summand $m$ acts trivially. Normal generation by $m$ makes the
entire image act trivially there.
The active summand is consequently faithful on the image.

A finite-order complex matrix is diagonalizable. Displacement rank
one therefore gives the spectrum $\mu,1,\ldots,1$ on $U$, with
$\mu\ne1$. Conjugates have the same spectrum. Normal generation
also implies $\langle C\rangle=W$.
\end{proof}

Put $r=\dim U$, and regard $W$ as an irreducible finite complex
reflection group on $U$. Let $\rr(W)$ be its minimum number of
reflection generators, with all reflection classes allowed.
We use the following universal class reduction, whose elementary
proof is retained as Lemma~\texttt{lem:class} of \cite{class}:
\begin{equation}\label{eq:class}
 \langle C\rangle=W
 \quad\Longrightarrow\quad
 \text{one hyperplane orbit},\qquad d_C(W)=\rr(W).
\end{equation}
More precisely, a member of $C$ generates the full cyclic pointwise
stabilizer of its hyperplane. Thus every reflection is a power
of a member of $C$ on the same hyperplane, which permits replacement
of any reflection generating set by members of $C$.

For clarity, the one-orbit step follows from orbit characters:
if $\mathcal O$ is a hyperplane orbit, the product of its linear
equations changes under $W$ by a character. That character is one
on a reflection whose hyperplane is outside $\mathcal O$, and is
its nontrivial normal eigenvalue on a reflection with hyperplane
in $\mathcal O$. A generating reflection class therefore misses
no orbit. Determinants then show its members have full
hyperplane-stabilizer order. These arguments apply to all ranks;
no special-knot result of \cite{class} is used here.

\begin{lemma}\label{lem:determinant}
Let $W$ be a finite complex reflection group, and let $L$ be the
least common multiple of the orders of all its reflections.
If one reflection conjugacy class generates $W$, then $W$
contains a reflection of order $L$.
\end{lemma}
\begin{proof}
The determinant of a reflection has exactly its order.
Since reflections generate $W$, the finite subgroup
$\det W\subset\C^*$ has order $L$.
If $C$ generates $W$, its common determinant generates $\det W$.
Every member of $C$ must therefore have order $L$.
\end{proof}

\section{The exact classification split}

Call an irreducible rank-$r$ reflection group well generated when
it has $r$ reflection generators. Cyclic rank one is included.
We import the complete irreducible classification and reflection
generator information from Shephard--Todd
\cite[Table~VII and its footnote, printed p.~301]{st54}.
Apart from the following cases, an irreducible group of rank $r$
is generated by $r$ reflections:
\[
 \begin{gathered}
 G(m,p,r)\quad(r\ge2,\ 1<p<m),\\
 G_7,G_{11},G_{12},G_{13},G_{15},G_{19},G_{22},\qquad G_{31}.
 \end{gathered}
\]
The seven numbered groups on the middle line have rank two;
$G_{31}$ has rank four.
All irreducible primitive rank-two groups are generated by at most
three reflections \cite[Section~4, printed p.~280]{st54}.
The classification, including completeness of this exception list,
is an imported theorem.

The mixed imprimitive groups have coordinate and off-coordinate
reflecting hyperplanes. Monomial transformations preserve the
distinction, so they have at least two hyperplane orbits
\cite[Section~2, printed pp.~276--277]{st54}.
They are excluded by \eqref{eq:class}.

Four of the rank-two exceptions are also excluded.
The complete reflection-order data in
\cite[Tables~I--III, printed pp.~281--282]{st54} give:
\[
\begin{array}{c|c|c}
 \text{group}&\text{orders of its reflections}&L\\\hline
 G_7&2,3&6\\
 G_{11}&2,3,4&12\\
 G_{15}&2,3&6\\
 G_{19}&2,3,5&30 .
\end{array}
\]
None contains a reflection of order $L$, contradicting
Lemma~\ref{lem:determinant}. The column convention for these tables
is stated on printed page~280: $r_q$ counts reflections of order $q$.

The remaining rank-two exceptions $G_{12},G_{13},G_{22}$ are
primitive and all their reflections have order two
\cite[Section~4 and Tables~II--III, printed pp.~279--282]{st54}.
By the rank-two generator bound and \eqref{eq:class},
\[
                         d_C(W)\le3
\]
in these cases. The next construction supplies a three-dimensional
irreducible KCH representation from their two-dimensional models.

\section{The determinant-twisted adjoint construction}

\begin{lemma}\label{lem:adjoint}
Let $W\subset\GL_2(\C)$ be finite and primitive, and suppose it is
generated by involutory reflections. On $\mathfrak{sl}_2(\C)$ define
\[
                     F(g)(X)=\det(g)\,gXg^{-1}.
\]
Then $F$ is an irreducible three-dimensional complex representation.
Every involutory reflection of $W$ acts in $F$ with eigenvalues
$-1,1,1$.
\end{lemma}
\begin{proof}
The determinant character and adjoint action are homomorphisms, so
$F$ is a representation. Since the generators have determinant
$-1$, the determinant character takes only the values $\pm1$.

Unitarize $W$. Conjugation then preserves the real three-dimensional
space $\mathcal H$ of traceless Hermitian matrices and its positive
inner product $\langle X,Y\rangle=\operatorname{tr}(XY)$.
Its complexification is $\mathfrak{sl}_2(\C)$.
The determinant twist is a real sign, so $F$ is real orthogonal on
$\mathcal H$. Multiplying operators by nonzero scalar characters
does not change their invariant subspaces.

Suppose the complex representation is reducible. A reducible
three-dimensional unitary representation has an invariant complex
line $L$. If $L=\overline L$, it contains a nonzero real vector,
giving an invariant real line in $\mathcal H$.
If $L\ne\overline L$, the conjugation-stable plane
$L+\overline L$ is the complexification of an invariant real
plane. Its real orthogonal complement is an invariant real line.
Thus in either case the adjoint action preserves a real line
$\mathbb R H$ with $H\in\mathcal H$ nonzero.

The Hermitian matrix $H$ has two distinct eigenvalues $t,-t$.
Norm preservation gives $gHg^{-1}=\pm H$ for every $g\in W$,
so $W$ permutes the two eigenspaces of $H$.
These are complex lines whose direct sum is $\C^2$.
That is an imprimitivity decomposition, contradicting primitivity.
Hence $F$ is irreducible.

Finally an involutory reflection is unitarily conjugate to
$\operatorname{diag}(-1,1)$. Its adjoint action fixes the traceless
diagonal direction and negates the two real off-diagonal
directions. Multiplication by its determinant $-1$ gives
$-1,1,1$ for $F$.
\end{proof}

The lemma requires no identification of the resulting finite
three-dimensional image. In particular no quaternion model,
projective-group recognition or numerical computation is a premise.

\section{A representation realizing the class bound}

\begin{lemma}\label{lem:capacity}
Let a group $G$ be normally generated by $m$, and let
$\rho:G\to\GL(V)$ be a finite-dimensional complex representation
with finite image $W$ and $\rank(\rho(m)-I)=1$.
Write $C$ for the conjugacy class of $\rho(m)$ in $W$.
Then $G$ has an irreducible finite-image complex representation
$\psi:G\to\GL(U')$ such that
\[
 \rank(\psi(m)-I)=1,\qquad \dim U'\ge d_C(W).
\]
The meridian image $\psi(m)$ is diagonalizable with one nontrivial
eigenvalue. No augmentation correspondence is needed for this assertion.
\end{lemma}
\begin{proof}
Lemma~\ref{lem:active} gives a faithful irreducible natural action of
$W$ on $U$, where $r=\dim U$, and $C$ generates $W$.
The classification split above applies.

For a well-generated group, \eqref{eq:class} gives $d_C(W)\le r$.
Use the natural action on $U$, composed with $\rho$.

For $G_{12},G_{13},G_{22}$, compose $\rho$ with $F$ from
Lemma~\ref{lem:adjoint}. Surjectivity onto $W$ makes this composite
irreducible, and its distinguished image has eigenvalues $-1,1,1$.
Here $d_C(W)\le3$, so dimension three suffices.

For $G_{31}$, the exact reflection-class quotient retained in
\cite[quotient theorem and class-rank corollary]{g31} is
\[
                       G_{31}\twoheadrightarrow S_6,
\]
taking all sixty reflections onto the transpositions, four to one.
That reflection set is one class, and its class-generating rank is
five. The standard representation of $S_6$ on
\[
                 \{(z_1,\ldots,z_6)\in\C^6:\sum_i z_i=0\}
\]
is irreducible of dimension five; a transposition has one eigenvalue
$-1$ and four eigenvalues one. Compose this representation with the
quotient and $\rho$. Its dimension is $d_C(G_{31})=5$.

The mixed imprimitive and four mixed-order rank-two exceptions
cannot occur under the generating-class hypothesis.
The split is exhaustive. All the constructed representations
factor through finite groups, so their distinguished images are
diagonalizable; the displayed spectra have displacement rank one.
\end{proof}

\begin{proof}[Proof of Theorem~\ref{thm:domination}]
Apply Lemma~\ref{lem:capacity} to the knot group.
The constructed irreducible representation is a KCH representation
of dimension $d\ge d_C(W)$, with distinguished eigenvalue different
from one. Cornwell's representation-to-augmentation construction
\cite[Proposition~2.12, printed p.~9]{cor17}, with rank as in
\cite[Definition~1.8, printed p.~3]{cor17}, supplies an ordinary
$U=1$ augmentation of rank $d$.
Thus $d_C(W)\le d\le\ar(K)$. No longitude restriction is added.
\end{proof}

Allowing inverse-oriented meridians does not alter either result.
A generator in $C^{-1}$ can be inverted without changing the subgroup
generated. Thus the minimum for $C\cup C^{-1}$ equals $d_C(W)$.

\begin{corollary}\label{cor:finitecapacity}
For every knot, the maximum reflection-class generating rank
among its finite complex rank-one-meridian images equals the
maximum dimension of its irreducible finite-image complex KCH
representations. Both maxima exist.
\end{corollary}
\begin{proof}
Lemma~\ref{lem:capacity} gives the first maximum as a lower bound
for the second. Conversely, let an irreducible finite-image KCH
representation have dimension $d$ and image $W$.
Choose $d_C(W)$ generators from its meridian class.
Their nonzero rank-one displacement directions span a nonzero
invariant space of dimension at most $d_C(W)$.
Irreducibility makes it the whole representation, so
$d\le d_C(W)$.

Both sets of integers are nonempty: the cyclic representation obtained
from knot abelianization by sending a meridian to $-1$ supplies one.
Images of any meridional generating set show
$d_C(W)\le\operatorname{mr}(K)$, using inverses when necessary.
The preceding span argument bounds the other set as well.
Thus both sets lie in the finite set
$\{1,\ldots,\operatorname{mr}(K)\}$.
\end{proof}

Theorem~\ref{thm:domination} bounds this common value by $\ar(K)$.
The equality in Corollary~\ref{cor:finitecapacity} concerns
finite-image KCH representations; equality with unrestricted
augmentation rank is not asserted.

\section{One application and scope}

\begin{corollary}\label{cor:knots}
Let $K_s=K_*\mathbin{\#}{}^s3_1$, $s\ge0$, be the family in
\cite{new17}. For every finite complex rank-one-meridian image
$W$ of its knot group,
\[
                     d_C(W)\le s+3<s+4=\operatorname{mr}(K_s)=b(K_s).
\]
For $s=0$ the ceiling three is attained by the explicit finite
reflection image in that companion. No attainment assertion for
the ceiling $s+3$ when $s>0$ is needed here.
\end{corollary}
\begin{proof}
The companion establishes $\ar(K_s)=s+3$ and the separate
meridional and bridge equalities, using its explicit rank-three
witness, the common-parameter trefoil construction and independent
finite-action bounds. Apply Theorem~\ref{thm:domination}.
Its finite three-dimensional seed image has class-generating
rank exactly three.
\end{proof}

The domination theorem itself has no dependency on this application
or on a particular knot census.
Its essential imported facts are the primary classification and
generator data, the universal class reduction, the explicit
$G_{31}$ quotient, and the ordinary augmentation construction.
The determinant test and twisted-adjoint irreducibility proof
are written above.
Exact source hashes and the table locators are retained in
\path{finite_complex_augmentation_domination_source_audit.json}.
The classification is not reproved here; no optional finite
computation is substituted for it.
This is a written comparison with declared mathematical imports,
not formal verification or external specialist review.
Worldwide priority has not been established.

\begin{thebibliography}{9}\small
\bibitem{st54}
G. C. Shephard and J. A. Todd,
\emph{Finite unitary reflection groups},
Canadian Journal of Mathematics \textbf{6} (1954), 274--304.
Sections~2 and~4, Tables~I--III, Table~VII and its footnote.
\url{https://doi.org/10.4153/CJM-1954-028-3}.
\bibitem{cor17}
C. R. Cornwell,
\emph{KCH representations, augmentations, and $A$-polynomials},
Journal of Symplectic Geometry \textbf{15} (2017), 983--1017.
Retained \href{https://arxiv.org/abs/1310.7526v3}{arXiv:1310.7526v3},
Proposition~2.12 and Definition~1.8.
\bibitem{class}
\emph{A finite complex reflection boundary for sixteen knots},
retained \path{FINITE_COMPLEX_REFLECTION_BOUNDARY.tex},
universal Lemma~\texttt{lem:class} only.
\bibitem{g31}
\emph{The symmetric-group quotient of the reflection class of $G_{31}$},
retained \path{G31_COXETER_QUOTIENT.tex},
Theorem~\texttt{thm:quotient} and
Corollary~\texttt{cor:classrank}.
\bibitem{new17}
\emph{A generic rank-one obstruction for a seventeen-crossing knot},
retained \path{NEW_SEVENTEEN_GENERIC_KCH.tex},
Corollaries~\texttt{cor:gap}, \texttt{cor:sums} and
\texttt{cor:finiteattainment}.
\end{thebibliography}
\end{document}
